\chapter{研究有效性、样本外与交易摩擦}
\label{ch:research-validity}

同一组历史数据，可以因模型选择、信息时点或成交假设不同而得到不同收益。本章以一个小型因子研究为例，把统计推断与交易中的具体数量联系起来，考察一个经验结论究竟说明了什么。

案例给出三个测试日的收益贡献 $0.02,-0.01,0.03$，以及从 $m=20$ 个候选假设中选出的原始 $p$ 值 $0.002$。下面将分别解释收益的计算、选择对显著性的影响，以及成本对结论的改变。

\section{研究对象、信息时点与收益定义}

计算收益之前，需要明确每个收益值的含义、形成信号时已经知道的信息，以及模型是否利用了评价区间中的数据。这些选择决定了收益曲线对应的统计问题。

在这个例子中，研究者于 2024 年 6 月 28 日结束模型选择，从 7 月 1 日开始观察测试区间；三个测试日各产生一笔按同一资本基数计量的收益贡献。先从最简单的收益与成本计算开始。

\begin{example}
三笔毛收益贡献之和为
\[
  R^{\mathrm{gross}}=0.02-0.01+0.03=0.04.
\]
每笔成本为 $0.002$，三笔共 $0.006$，所以
\[
  R^{\mathrm{net}}=R^{\mathrm{gross}}-C=0.04-0.006=0.034.
\]
这里的加法成立，是因为三项被定义为同一资本基数上的收益贡献。若它们是连续持有期收益，应使用复利；若仓位不同，应先乘权重。定义错了，算术再精确也没有意义。
\end{example}

\MFQLead{时点数据：一行数据需要四个时间}

设研究者在决策时刻 $t$ 可使用的信息集合为 $\mathcal I_t$。一个特征 $X$ 可进入该决策，必须能证明 $X$ 对 $\mathcal I_t$ 可测；“数据库里现在有这一列”不是证明。对每一行至少记录：
\begin{itemize}
  \item \textbf{事件时间 $t_e$.} 经济事件实际发生的时间；
  \item \textbf{可得时间 $t_a$.} 市场参与者按指定数据源首次能够获得该值的时间；
  \item \textbf{决策时间 $t_d$.} 策略确定特征并形成订单意图的时间；
  \item \textbf{目标起点 $t_y$.} 被预测收益开始累积的时间。
\end{itemize}
本例时间均为同一市场本地时间。历史数值应对应当时发布的版本；若 $t_e$ 表示已发生事件，常有 $t_e\leq t_a$，而决策合法性直接要求
\[
 t_a\leq t_d<t_y.
\]
预测性公告可能早于所预测事件，不能把 $t_e\leq t_a$ 当成所有字段的普遍公理。

若日报表在 7:30 对应前一交易日事件、8:00 才可获得，而决策在 9:25、收益从 9:30 起算，则该数据在决策时可用。把可得时间改成 9:31，即使事件发生在早晨，也构成前视。

\MFQLead{泄漏不只来自“用了未来价格”}

常见泄漏至少有五类：
\begin{enumerate}[leftmargin=*]
  \item \textbf{特征泄漏：}使用决策后才发布或修订的数据；
  \item \textbf{标签泄漏：}标准化、缺失填补或特征筛选看到了测试期标签；
  \item \textbf{横截面泄漏：}在当日收盘后才能完整知道的全市场统计被用于盘中决策；
  \item \textbf{资产池泄漏：}用今天仍存续的证券回填历史，遗漏退市、停牌或不可交易资产；
  \item \textbf{选择泄漏：}反复观察测试集，再把最后一次配置称为“样本外”。
\end{enumerate}
因此 point-in-time（时点）数据不仅要求值的历史版本，还要求资产池、日历、复权因子、分类标签与数据发布时间都可回到当时状态。最终数据库若只保留“最新修订值”，便无法恢复当时的信息集。

\begin{example}
财报期末是 $t_e$，公告首次发布时间才是 $t_a$。把财报期末日期当作可得日期，会让研究者提前数周或数月知道结果。正确做法是保存公告版本与发布时间，并对同日公告明确决策是在公告前还是公告后。
\end{example}

\section{样本外与选择自由度}

\begin{definition}[样本外]
样本外（out of sample）区间是从未参与模型拟合、特征选择、超参数选择、停止规则或研究叙事选择的数据。只要看过该区间并据此修改研究，它就不再是最终样本外。
\end{definition}

% Generated by tools/build_figures.py; edit figures/figure-specs.json instead.
\MFQFigure{tex/figures/generated/upper-ch16-purge-embargo.pdf}{训练标签与验证区间重叠时，purge 删除相应训练样本；embargo 在需要时增加时间缓冲，长度取决于标签及特征的信息范围。}{upper-ch16-purge-embargo}

% Short alias: \begin{figure}[htbp]
  \centering
  \begin{tikzpicture}[x=0.88cm,y=0.72cm,>=stealth]
    \draw[->] (0,0) -- (12.2,0) node[right] {时间};
    \fill[blue!20] (0.5,0.35) rectangle (4.6,1.05);
    \fill[red!18] (4.6,0.35) rectangle (6.0,1.05);
    \fill[green!22] (6.0,0.35) rectangle (8.8,1.05);
    \fill[orange!22] (8.8,0.35) rectangle (9.8,1.05);
    \fill[blue!20] (9.8,0.35) rectangle (11.7,1.05);
    \node at (2.55,0.70) {训练};
    \node at (5.30,0.70) {purge};
    \node at (7.40,0.70) {验证};
    \node at (9.30,0.70) {embargo};
    \node at (10.75,0.70) {训练};
    \draw[<->,thick] (3.4,1.45) -- (6.8,1.45) node[midway,above] {标签窗口重叠};
  \end{tikzpicture}
  \caption{purge 删除与验证标签窗口重叠的训练样本；embargo 隔离验证后的慢变信息。}
\end{figure}


最简单的时间切分为训练集、验证集和一次性测试集：训练集估计参数，验证集选择模型，最终测试集只在研究设计确定后打开一次。金融分布会变化，因此开发阶段更常使用滚动向前验证（walk-forward validation）：在每个时点只用过去窗口拟合，再预测紧随其后的未来窗口。每一折都保持时间方向，但所有折仍属于模型开发证据；不能把反复优化后的平均得分当作未见过的最终测试。

\MFQLead{重叠标签、purge 与 embargo}

若特征时点为 $t$，标签使用 $[t,t+h]$ 的未来收益，相邻样本的标签区间可能重叠。普通随机交叉验证会同时把相互共享未来价格的样本放进训练和验证。解决思路是：
\begin{itemize}[leftmargin=*]
  \item 从训练折删除标签窗口与验证折重叠的样本（purge）；
  \item 在验证窗口附近留出缓冲区（embargo），防止慢变特征或执行状态跨边界；
  \item 若超参数也被选择，在外层评估、内层选择，避免同一折同时承担选择与最终评分。
\end{itemize}


\begin{heuristic}
purge 长度与 embargo 长度不是通用常数，它们应由标签跨度、特征记忆和研究机制决定。机械地套一个百分比并不能证明独立性；这里给出的是研究设计启发，不是保证样本独立的定理。
\end{heuristic}

\begin{example}[一张可复核的时间窗口表]
明确把第 $t$ 日标签定义为第 $t$ 至 $t+4$ 日的收益贡献，训练候选为 1--20 日，验证为 21--25 日。验证标签的并集为第 21--29 日；训练中第 17--20 日标签与该并集相交，所以删除后留下 1--16 日。若改用 $(t,t+5]$，必须重新按端点计算，不能机械照搬行号。

过去三日的滚动特征在第 21 日使用第 18--20 日数据，本身符合向前预测的信息要求，并不自动要求再删三天。是否另加缓冲应由延迟、标签完成时间和欲控制的依赖决定。purge 删除共享标签信息，不能保证所有统计依赖已经消失。

\end{example}

\MFQLead{分布漂移改变结论的适用域}

设训练分布为 $P_{\mathrm{train}}(X,Y)$，部署分布为 $P_{\mathrm{live}}(X,Y)$。样本外测试估计的是某个历史区间的风险，不自动保证未来风险相同。至少要区分协变量漂移、条件关系漂移、资产池变化和交易制度变化。若测试期恰好只覆盖单一市场状态，结论应限定到该状态，而不是写成“策略稳定有效”。

\MFQLead{多重检验：研究次数也会改变证据}

若在每个真零假设下用显著性水平 $\alpha$ 独立检验 $m$ 次，至少出现一次假阳性的概率为
\[
  1-(1-\alpha)^m.
\]
当 $\alpha=0.05,m=20$ 时，该概率约为 $64.15\%$。即使检验并不独立，研究者也不能因为只展示赢家而把搜索过程从证据中删除。

Bonferroni 方法用并集界控制族错误率（family-wise error rate）：
\[
  \Pr\left(\bigcup_{j=1}^m\{p_j\le\alpha/m\}\right)
  \le \sum_{j=1}^m \Pr(p_j\le\alpha/m)\le\alpha.
\]
因此本章案例的单项阈值为
\[
  \frac{0.05}{20}=0.0025,
\]
最终 $p=0.002$ 通过；若改成 $0.003$，普通 $5\%$ 检验仍会拒绝，但在当前错误率标准下应判为不支持。Bonferroni 不要求检验独立，但可能保守；当目标是控制发现集合中的错误比例时，可考虑第 10 章的 Benjamini--Hochberg 程序。选择哪种错误率，应在看结果前说明。

\MFQLead{什么算一次尝试}

研究自由度（research degrees of freedom）包括：信号定义、滞后、窗口、缺失值规则、资产池、行业中性化、成本模型、持有期、评价指标和停止时间。只把最终 notebook 中的 20 列计为 $m$，却忽略删除过的几百个版本，仍会低估选择偏差。更透明的做法是逐次记录假设族、数据版本、时间、主要指标和放弃原因。预注册不是禁止探索，而是把探索性结论与确认性检验分开。

\section{交易摩擦、市场条件与净收益}

最小成本模型可写为
\[
  R^{\mathrm{net}}
  =R^{\mathrm{gross}}
  -C_{\mathrm{commission}}
  -C_{\mathrm{spread}}
  -C_{\mathrm{impact}}
  -C_{\mathrm{borrow}}
  -C_{\mathrm{other}}.
\]
佣金常与成交金额近似成比例；买卖价差取决于订单方向与报价；市场冲击随参与率、订单规模、波动率和执行时长变化；借券成本与可得性对空头策略构成状态约束。把所有成本压成一个固定基点，只适合明确标记的教学近似。

\MFQLead{换手、冲击与容量}

若权重从 $w_{t-1}$ 调整到 $w_t$，一种双边换手定义为
\[
  \operatorname{TO}_t=\sum_i |w_{i,t}-w_{i,t-1}|.
\]
也有文献使用其一半，因此比较换手率时需使用同一约定。线性成本近似可写为 $c\operatorname{TO}_t$，但冲击通常不是线性的。策略容量（capacity）是净绩效在资金规模、成交量约束与冲击模型下仍满足目标的最大资金规模；它不是“回测成交金额的平均值”。

实现差额（implementation shortfall）比较决策时基准价格与实际执行结果，把延迟、价差、冲击和未成交机会成本汇总到同一执行基准。只用日收盘价成交会遗漏这些路径。一个可信回测至少要报告：订单生成时刻、可用报价、成交规则、未成交处理、成交量限制和成本参数来源。

\begin{example}
若毛收益贡献为 4\%，逐项成本记录扣除 0.6\% 后剩 3.4\%。若把资金规模扩大十倍，而冲击随参与率非线性增长，就不能把 3.4\% 原样外推。容量问题需要重新计算成交与冲击，而不是只把盈亏乘十。
\end{example}

\MFQLead{把容量写成函数而不是一句话}

设单日成交量为 $V_t$、策略下单量为 $q_t$，参与率为
$u_t=|q_t|/V_t$，一个简单的冲击模型可写成
\[
 C_{\rm impact}(q_t)
 =\eta\,\sigma_t\,p_tV_t\,u_t^{\,\gamma},
 \qquad \gamma>1.
\]
其中 $V_t,q_t$ 以股为单位，$p_t$ 为每股价格，$\sigma_t$ 是无量纲收益波动率，所以 $C_{\rm impact}$ 是金额；换成收益成本还需除以资本。当资金规模放大 $s$ 倍而信号权重不变时，$u_t$ 也近似放大 $s$，成本按 $s^\gamma$ 而非 $s$ 增长。容量应定义为满足
$\mathbb E[R^{\rm net}(s)]\ge r_{\min}$、成交约束和风险约束的最大 $s$；只把回测收益乘以 $s$ 是把线性收益误当成线性可成交性。

\MFQLead{全球基线与 A 股条件}

全球主线不假设某个交易所的固定参数。所有市场都要核验交易日历、点时资产池、可成交报价、费用、价差、冲击、借券与公司行动。A 股研究在相关资产和研究期内还常需显式加入：买入证券的库存可售时间约束、不同板块与证券状态的涨跌幅及有效申报价格、停复牌、point-in-time 复权与公司行动、融资融券标的和券源资格。这些是成交与样本选择条件，不是数学定理。

规则会变化，也因板块、证券状态和交易方式不同而变化。本书因此不把某个涨跌幅百分比写进通用公式；复现报告应固定研究期间适用的规则版本。作为当前规则核验入口，可查上海证券交易所《交易规则（2026 年修订）》\footnote{\url{https://www.sse.com.cn/lawandrules/sselawsrules2025/stocks/exchange/c/c_20260424_10816482.shtml}}与深圳证券交易所《交易规则（2026 年修订）》\footnote{\url{https://docs.static.szse.cn/www/lawrules/rule/trade/current/W020260424690713155663.pdf}}，再根据研究期回溯当时生效版本。

\begin{center}
\begin{tabular}{p{0.24\textwidth}p{0.31\textwidth}p{0.34\textwidth}}
\toprule
审计对象 & 全球通用基线 & A 股条件化问题 \\
\midrule
可交易日与资产池 & 交易所日历、上市状态 & 停复牌、板块和风险状态的历史版本 \\
订单与成交 & 报价、方向、数量、延迟 & 涨跌幅与有效申报范围是否使成交可行 \\
库存与空头 & 库存、借券可得性与成本 & 买入后可售约束、融券标的与券源资格 \\
价格与公司行动 & point-in-time 价格口径 & 复权因子的公布时间及除权除息处理 \\
\bottomrule
\end{tabular}
\end{center}

一个 point-in-time 例子能看出这些条件为何不能省略：若 6 月 28 日收盘后形成信号，7 月 1 日开盘执行，若信号在 6 月 28 日收盘后立即固定，6 月 30 日和 7 月 1 日 9:05 才发布的字段都不能回填该信号。若事先允许开盘前重算，则应把新的决策时刻和数据截止时刻写清；“开盘成交”本身没有定义信号截止。若数据库只保存“报告期”而不保存“可得时间”，回测即使使用正确的数值，也无法证明当时可得。因此，解释历史决策还需要对应当时的数据版本、发布日期、时区及公司行动调整方式。


\noindent\textbf{延伸阅读。}数据窥探与回测过拟合的正式处理参见\cite{white2000reality,bailey2017backtest}。

\section*{章末练习}
\begingroup\small
\subsection*{口述概念}
\begin{enumerate}[nosep,leftmargin=*]
  \item 区分事件时间、可得时间、决策时间与目标起点；哪两个不等式直接阻止前视？
  \item 为什么反复查看一个“测试集”后，它不再是最终样本外？
  \item 复现、统计有效和可交易三者分别回答什么问题？为何互不替代？
\end{enumerate}

\subsection*{笔试推导}
\begin{enumerate}[nosep,leftmargin=*]
  \item 在独立真零假设下推导 $m$ 次水平 $\alpha$ 检验至少一次假阳性的概率，并计算 $m=20,\alpha=0.05$。
  \item 用并集界证明 Bonferroni 阈值 $\alpha/m$ 控制族错误率；判断 $p=0.002$ 与 $p=0.003$ 是否通过本章协议。
  \item 三笔等资本收益贡献为 $0.02,-0.01,0.03$，每笔佣金、价差、冲击为 $0.0005,0.0007,0.0008$。推导毛收益、总成本和净收益，并说明何时不能直接相加。
\end{enumerate}

\subsection*{数值编程}
\begin{enumerate}[leftmargin=*]
  \item 用闭区间标签 $[t,t+4]$ 算出训练 1--20 日与验证 21--25 日的重叠。
  \item 画出毛贡献 $0.003$、换手 2、其他成本 $0.0004$ 时净收益对成本率的曲线。
  \item 构造可得时刻晚于决策的例子，说明如何改变模型输入。
\end{enumerate}

\subsection*{研究判断}
\begin{enumerate}[nosep,leftmargin=*]
  \item 研究者尝试 500 个因子，只报告最佳者的 $p=0.01$ 和一次 bootstrap 区间。列出至少四个缺失证据，并设计修复协议。
  \item 美股日频策略直接迁移到 A 股，假定每日收盘价全部成交且可立即卖出。哪些结论可能被制度边界改变？哪些数学结论仍是通用的？
  \item 一个策略成本前 Sharpe 很高，成本后仍为正，但资金扩大十倍后仍沿用固定 5bp 成本。你会拒绝哪项结论，要求哪些容量证据？
\end{enumerate}
\endgroup

\subsection*{基础衔接：独立计算}
每期毛收益 $g=0.003$，双边换手 $\tau=2$，除线性成本外还有收益成本 $h=0.0004$。求盈亏平衡的单位换手成本。

\subsection*{分级提示}
\begingroup\small
\begin{enumerate}[nosep,leftmargin=*]
  \item \textbf{口述概念：}先画信息集与研究者可选择动作；复现只约束输入到输出，统计有效约束推断，可交易性还约束订单到成交。
  \item \textbf{笔试推导：}独立时先算“一次假阳性都没有”；Bonferroni 使用并集界而非独立性；收益记录要先声明同一资本基数。
  \item \textbf{数值编程：}先写出本章公式中的目标量，再显示中间计算。改变参数前预测结果方向，运行后说明差异来自模型、抽样还是数值近似。
  \item \textbf{研究判断：}把选择过程当作算法的一部分；A 股特例进入成交与资产池协议；容量要求冲击随规模或参与率变化。
\end{enumerate}
\endgroup


\MFQLead{用一张研究审计卡收束证据}

章末请整理一张研究审计卡：目标参数与适用市场、四个时间字段、样本外边界、研究次数、毛收益到净收益的成本分解、成交和容量假设、数值容差，以及一个泄漏反例和一个选择偏差反例。这张卡会成为上册综合项目复核研究主张的共同起点；下一章会把它放进一个完整案例，继续追问这些证据能够支持什么结论。

\noindent\textbf{本章要点。}经验结果与信息时点、模型选择和成交假设相联系。样本外评价用于考察选定方法在新数据上的表现，成本与容量分析则连接统计预测和实际收益。三笔交易可以说明一种计算关系，长期平均收益与未来表现仍需要相应的数据和推断。